\documentclass[10pt,a4paper]{amsart}
\usepackage[margin=19mm]{geometry}
\usepackage{amsmath,amssymb,mathtools}
\usepackage[scr=rsfs,cal=euler]{mathalfa}
\usepackage{xfrac}
\usepackage[unicode,hidelinks,bookmarks=false]{hyperref}
\newcommand{\ie}{i.e.\ }
\newcommand{\cf}{cf.\ }
\newcommand{\resp}{respectively\ }
\newcommand{\Subsection}{\subsection}
\providecommand{\nobreakdash}{\nobreak\hbox{-}}
\setlength{\emergencystretch}{2em}
\multlinegap0pt
\usepackage[shortlabels]{enumitem}
\usepackage{amssymb}
\usepackage{mathtools}
\usepackage{tikz}
\usepackage{xcolor}
\newtheorem{theorem}{Theorem}
\newtheorem{proposition}{Proposition}
\newcommand{\Ad}{\operatorname{Ad}}
\newcommand{\R}{\mathbb R}
\newcommand{\Ric}{\operatorname{Ric}}
\newcommand{\dd}{\mathrm d}
\newcommand{\eps}{\varepsilon}
\title{Future Completeness, $C^0$-Inextendibility and Cauchy Horizons\\ in Homogeneous Einstein Spacetimes}
\author{Bobby Eka Gunara}
\date{Source: arXiv:2609.11979v1 (CC BY 4.0)}
\begin{document}
\maketitle

\noindent\emph{Source and licence.} Future Completeness, $C^0$-Inextendibility and Cauchy Horizons\\ in Homogeneous Einstein Spacetimes. Version arXiv:2609.11979v1 (CC BY 4.0), distributed by arXiv under the Creative Commons Attribution 4.0 International licence, http://creativecommons.org/licenses/by/4.0/, as recorded in the arXiv licence metadata for this exact version and shown on the abstract page https://arxiv.org/abs/2609.11979v1. Full licence notice: \texttt{licenses/arxiv-2609.11979-CC-BY-4.0.txt}.\par\smallskip

\noindent\textit{Editorial scope.} Counted unit: the article's proposition prop:einstein-optimal-separation (source0.tex lines 3663--3710). The article's complete original proof of it is delivered with it (source0.tex lines 3712--3913, one \texttt{proof} environment of 202 lines). No mathematics is written, repaired or re-derived: the statement and the proof are the article's own lines, in the article's own notation, and no retained line is substituted or re-flowed.  Retained uncounted as the source-local context the proof needs: lines 291--297; lines 300--353; lines 368--373; lines 1887--1892; lines 2308--2313; lines 2340--2363; lines 5196--5202; lines 6812--6858; lines 6951--6954; lines 6957--6975.  Nothing was omitted from any retained span, so the package carries no omission record.  No retained line contains a citation, so the wrapper carries no bibliography and no bibliography entry.  No retained line contains a figure environment, an image-inclusion command or drawing code, so no image is delivered and the package declares no asset dependency.  Editorial formatting only, in addition: an AMS \texttt{amsart} preamble that restates the article's own declarations and macros used by the retained lines, namely the environments \texttt{theorem}, \texttt{proposition} on separate counters, together with the standard packages those lines need.  The retained environments print in this document as Theorem 1, Proposition 1 in order of appearance, whereas the article prints the same items with the numbering of its own full document.  The complete native source of the article as distributed in the arXiv e-print for this exact version is bundled unchanged as \texttt{sources/210066/source0.tex}.
\begin{equation}
 \Ric(g)=\Phi^*\mathsf G+\frac{2\mathcal V(\Phi)}{n-1}g,
 \qquad D_t\dot\Phi+\theta\dot\Phi+
             \operatorname{grad}_{\mathsf G}\mathcal V=0,
 \qquad \theta=\tfrac12\operatorname{tr}_{h_t}\dot h_t,
 \label{eq:main-sigma-normalization}
\end{equation}

\begin{theorem}[Future completeness and the G\"odeke-Rendall conjecture]
\label{thm:homogeneous-future}
Let $X=L/H$ be a connected homogeneous space of dimension $n\ge2$,
where $L$ is a connected Lie group and $H$ is closed.  Let
$(\mathcal T,\mathsf G)$ be a finite-dimensional connected complete
Riemannian manifold and let $\mathcal V\in C^\infty(\mathcal T)$.
Suppose that $g=-\dd t^2+h_t$, with $h_t$
$L$-invariant, and $\Phi=\Phi(t)$ form a maximal homogeneous
solution of~\eqref{eq:main-sigma-normalization}.
Assume that $\theta(t_0)=\theta_0>0$ and that
$R(h_t)\le0$ and $\mathcal V(\Phi(t))\ge0$ throughout the future
interval.  Then the future
proper-time endpoint is infinite.  The spacetime is future timelike
and null geodesically complete on $X$, and on every quotient
$\Gamma\backslash X$ for which $\Gamma<L$ acts freely and properly
discontinuously, in each of the following cases:
\begin{enumerate}[label=\textup{(\roman*)}]
\item $\mathcal V\equiv0$ and $2\le n\le9$;
\item $2\le n\le5$;
\item $n\ge3$ and
\[
 \int_{t_0}^{\infty}
       \mathcal V(\Phi(t))^{\beta_n/2}\,\dd t=\infty,
 \qquad \beta_n=\frac{(n-1)^2}{4n}.
\]
\end{enumerate}
For any future causal geodesic with affine parameter $\lambda$,
\[
 E(t):=\frac{\dd t}{\dd\lambda}
 \le E(t_0)\left(\frac{\theta_0}{\theta(t)}\right)^{\beta_n},
 \qquad
 \beta_2=0,\quad \beta_n=\frac{(n-1)^2}{4n}\quad(n\ge3),
\]
and
\[
 \lambda(t)-\lambda(t_0)\ge
 \frac{1}{E(t_0)\theta_0^{\beta_n}}
 \int_{t_0}^{t}\theta(s)^{\beta_n}\,\dd s.
\]
Thus divergence of the last integral is an all-dimensional completeness
criterion.  When $\mathcal V\equiv0$, one may replace $\beta_n$ by
$c_2=0$, $c_3=1/3$ and $c_n=(\sqrt n-1)/2$ for $n\ge4$.
If $\mathcal V(\Phi(t))\ge V_*>0$, then
$\theta(t)\ge\sqrt{2nV_*/(n-1)}$, so each geodesic has uniformly
bounded energy in every dimension.  In particular, for a connected,
simply connected four-dimensional Lie group $G\not\simeq
\mathrm{SU}(2)\times\mathbb R$, every homogeneous vacuum development on
$G$ which is expanding at one slice is future timelike and null
geodesically complete, as conjectured by G\"odeke and Rendall.
The real-scalar equations
\eqref{eq:main-einstein-normalization} correspond to
$(\mathcal T,\mathsf G)=(\R,\dd\phi^2)$; a constant map with
$\mathcal V=0$ gives vacuum.
\end{theorem}

\begin{equation}
 \Ric(g)=\dd\phi\otimes\dd\phi+
 \frac{2}{n-1}\mathcal V(\phi)g,
 \qquad \Box_g\phi=\mathcal V'(\phi)
 \label{eq:main-einstein-normalization}
\end{equation}

\begin{equation}\label{eq:kappa-opt}
 k_I(\xi)=\sup_{0\ne\lambda\in Z^*}
 \frac{|\lambda(Q_I(\xi))|}
 {\displaystyle\int_I f_\xi(t)
          \sqrt{A_t^{-1}(\lambda,\lambda)}\,\dd t}.
\end{equation}

\begin{equation}
 h_{t_j}^{-1}(\alpha,\alpha)\longrightarrow0
 \quad\text{for some fixed }0\ne\alpha\in
 \operatorname{Ann}[\mathfrak n,\mathfrak n],\quad t_j\to t_-.
 \label{eq:fixed-closed-growth}
\end{equation}

\begin{theorem}[Continuous inextendibility on two-step quotients]
\label{thm:two-step-obstruction}
Suppose $\kappa_{\rm opt}(T)\le\tfrac12$ and
\begin{equation}
 \int_{t_-}^{T}\frac{N(t)}{m_h(t)}\,\dd t<\infty.
 \label{eq:schur-cone-integral}
\end{equation}
The spacetime on $\mathsf N$, and on $\Gamma\backslash\mathsf N$ for
every discrete subgroup $\Gamma<\mathsf N$, is past
$C^0$-inextendible if at least one of the following holds as $t\to t_-$:
\begin{enumerate}[label=\textup{(\roman*)}]
\item $\limsup\lambda_{\max}(\overline h_t)=\infty$;
\item $\limsup\sqrt{\det H_t\det A_t}=\infty$.
\end{enumerate}
The first condition includes horizontal growth and the fixed-covector
test \eqref{eq:fixed-closed-growth}.  No rationality condition relative
to the discrete subgroup is required.  Future timelike geodesic
completeness implies global $C^0$-inextendibility.  The stronger condition
$\int_{t_-}^T N/m_{\rm Sch}<\infty$ is sufficient for the cone hypothesis.
With a strict margin in the filling bound, the conclusion is stable under
the two-sided metric comparisons of Proposition~\ref{thm:stability}; it
also extends to invariant shifts under the hypotheses of
Proposition~\ref{prop:shift}.
\end{theorem}

\begin{equation}
 \kappa_{\rm cen}^0(T)
 \le\frac12\left(\sum_{\alpha=1}^r
       \lambda_{\max}(S_\alpha(T))^2\right)^{1/2}
 \longrightarrow0\qquad(T\downarrow0).
 \label{eq:anisotropic-bracket-matrix-bound}
\end{equation}

\begin{proposition}[Quantitative two-sided comparison]
\label{thm:stability}
Let
\[
 g=-N^2\dd t^2+h_t,\qquad
 \widetilde g=-\widetilde N^2\dd t^2+\widetilde h_t
\]
be invariant orthogonal evolutions on the same two-step group, or data
descending to the same quotient.  Suppose that, on $(t_-,T_0)$,
\begin{equation}
 (1-\eps)N\le\widetilde N\le(1+\eps)N,\qquad
 (1-\eps)h_t\le\widetilde h_t\le(1+\eps)h_t,\qquad 0<\eps<1.
 \label{eq:metric-comparison}
\end{equation}
Then, for $T\le T_0$ and
$C_\eps=(1+\eps)^{3/2}/(1-\eps)$,
\begin{equation}
 \widetilde\kappa_{\rm opt}(T)
 \le C_\eps\,\kappa_{\rm opt}(T),\qquad
 \widetilde\kappa_{\rm an}(T)
 \le C_\eps\,\kappa_{\rm an}(T).
 \label{eq:kappa-stability}
\end{equation}
If $\kappa_{\rm opt}(T)\le\tfrac12-\delta$ for
$0<\delta<\tfrac12$ and
\begin{equation}
 \frac{(1+\eps)^{3/2}}{1-\eps}
 \left(\frac12-\delta\right)\le\frac12,
 \label{eq:admissible-epsilon}
\end{equation}
then $\widetilde\kappa_{\rm opt}(T)\le\tfrac12$.
The largest admissible $\eps=\eps_*(\delta)$ is the unique root in
$(0,1)$ of equality in \eqref{eq:admissible-epsilon}, and
\begin{equation}
 \eps_*(\delta)=\frac45\delta+O(\delta^2)
 \qquad(\delta\downarrow0).
 \label{eq:epsilon-expansion}
\end{equation}
The comparison also preserves finiteness of $\int_{t_-}^T N/m_h$ and
each of the metric-divergence conditions in
Theorem~\ref{thm:two-step-obstruction}.
Hence every strict instance of that theorem remains past
$C^0$-inextendible under sufficiently small comparisons of the form
\eqref{eq:metric-comparison}.
The global conclusion also persists when the perturbed spacetime is
independently future timelike geodesically complete.
\end{proposition}

\begin{equation}
 g_\beta=-N^2\dd t^2+h_t(\omega+\beta\dd t,\omega+\beta\dd t)
 \label{eq:shift-metric}
\end{equation}

\begin{proposition}[Removal of invariant shift]
\label{prop:shift}
Let $a(t)$ solve $a^{-1}a'=-\beta(t)$.
The diffeomorphism $F(t,[x])=(t,[xa(t)])$ transforms
\eqref{eq:shift-metric} into
\[
 F^*g_\beta=-N^2\dd t^2+\bar h_t,\qquad
 \bar h_t=(\Ad_{a(t)^{-1}})^*h_t.
\]
The forms $H_t,A_t$ and the constants
$\kappa_{\rm opt},\kappa_{\rm an}$ are unchanged.
The conclusions of Theorem~\ref{thm:two-step-obstruction} therefore
hold with the cone integral computed using $m_{\bar h}$.
If the horizontal component $v_a(t)$ of $a(t)$ is bounded on the
past tail, finiteness of $\int_{t_-}^T N/m_h$ already implies that
cone condition.
It suffices that the horizontal drift $\beta_V$ be integrable;
no integrability of the central drift is needed.
\end{proposition}

\begin{proposition}[Distinct central scales in vacuum and scalar cosmologies]
\label{prop:einstein-optimal-separation}
Let $\mathsf N=H_3\times H_3$ and choose an invariant coframe
$\omega^1,\ldots,\omega^6$ satisfying
\[
 \dd\omega^1=-\omega^2\wedge\omega^3,\qquad
 \dd\omega^4=-\omega^5\wedge\omega^6,
\]
with the other coframe elements closed.  There exists a smooth
diagonal Einstein-massless-scalar solution
\[
 g=-\dd t^2+\sum_{i=1}^6 a_i(t)^2(\omega^i)^2,
 \qquad \operatorname{Ric}_g=\dd\phi\otimes\dd\phi,
\]
whose past asymptotics are
\begin{equation}
 \begin{split}
 a_i(t)&=t^{p_i}\bigl(1+O(t^{3/5})\bigr),\\
 \phi(t)&=q\log t+O(t^{3/5}),\\
 (p_1,\ldots,p_6)&=\frac1{20}(-8,-1,7,8,7,7),
 \qquad q=\frac{\sqrt{31}}{10}.
 \end{split}
 \label{eq:einstein-optimal-separation-data}
\end{equation}
The estimates hold after two applications of $t\partial_t$.
There is also a vacuum solution with the same error order and
Kasner data
\begin{equation}
 p^{\rm vac}=\frac1{20}
       (-8,3-\sqrt{78},3+\sqrt{78},8,7,7),\qquad q=0.
 \label{eq:vacuum-optimal-separation-data}
\end{equation}
For either maximal expanding homogeneous development,
\begin{align}
 cT^{3/10}\le\kappa_{\rm opt}(T)&\le C T^{3/10}
       &&\text{for sufficiently small }T>0,\notag\\
 \kappa_{\rm an}(T)&=\infty
       &&\text{for every }T\text{ in the proper-time domain}.
 \label{eq:einstein-optimal-separation}
\end{align}
Both developments are future timelike and null complete and globally
$C^0$-inextendible on $\mathsf N$ and every discrete left quotient.
The strict separation of the filling conditions, completeness and
inextendibility persist for a relative open neighborhood of either
displayed Kasner datum within the corresponding diagonal asymptotic
family.  The positive error and decay exponents may vary in these
neighborhoods.
\end{proposition}

\begin{proof}
Write $a_i=e^{\alpha_i}$ and use harmonic time $\tau$ with
$\dd t=e^S\dd\tau$, $S=\sum_i\alpha_i$.  Put
\[
 v_1=(-1,1,1,0,0,0),\qquad
 v_2=(0,0,0,-1,1,1),\qquad
 w_r=\mathbf1-v_r,
\]
and
\[
 \mathcal B_r=\frac12 e^{2w_r\cdot\alpha},\qquad r=1,2.
\]
For a diagonal invariant metric, the only nonzero orthonormal
spatial brackets have lengths
$e^{\alpha_1-\alpha_2-\alpha_3}$ and
$e^{\alpha_4-\alpha_5-\alpha_6}$.  In each block, the spatial Ricci
eigenvalues are one half the squared bracket length in the central
direction and minus one half in the two horizontal directions.
Independent sign changes of the four horizontal generators induce
the six distinct nontrivial characters
$\eta_1\eta_2,\eta_1,\eta_2,\eta_3\eta_4,\eta_3,\eta_4$
on the invariant frame.  They preserve the spacetime metric and
force all off-diagonal spatial and normal-spatial Ricci components
to vanish.  Thus the momentum constraint vanishes, and the
Einstein-scalar equations are
\begin{equation}
 \alpha''=\mathcal B_1v_1+\mathcal B_2v_2,
 \qquad \phi''=0,
 \qquad
 (S')^2-\sum_i(\alpha_i')^2
       = (\phi')^2+\mathcal B_1+\mathcal B_2.
 \label{eq:einstein-product-heisenberg-ode}
\end{equation}
Here and in the construction below primes denote $\tau$ derivatives.
The two curvature terms are coupled:
$w_r\cdot v_r=-2$ and $w_r\cdot v_s=1$ for $r\ne s$.
No decoupled formula for the Heisenberg family is used.

The data in \eqref{eq:einstein-optimal-separation-data} satisfy
\[
 \sum_i p_i=1,\qquad \sum_i p_i^2+q^2=1,\qquad
 \chi_1:=w_1\cdot p=\frac3{10},\qquad
 \chi_2:=w_2\cdot p=\frac7{10}.
\]
Set $\alpha=p\tau+u$ and $\phi=q\tau$.  On a sufficiently early
interval $(-\infty,\tau_0]$, solve the Volterra equation
\begin{equation}
 u(\tau)=\frac12\sum_{r=1}^2v_r
   \int_{-\infty}^{\tau}(\tau-s)
     e^{2\chi_rs+2w_r\cdot u(s)}\,\dd s.
 \label{eq:einstein-product-heisenberg-volterra}
\end{equation}
For completeness, take $\mu=3/5$ and the norm
$\|u\|_\mu=\sup_{\tau\le\tau_0}e^{-\mu\tau}|u(\tau)|$.
On a fixed ball in this space, first choose its radius larger than
the norm of the right side at $u=0$, and then decrease $\tau_0$.
The inequalities $2\chi_r\ge\mu$ and
\[
 \int_{-\infty}^{\tau}(\tau-s)e^{\mu s}\,\dd s
       =\frac{e^{\mu\tau}}{\mu^2}
\]
show that the right side preserves this ball.  The mean-value theorem
for the exponential gives a Lipschitz constant at most
$C e^{\mu\tau_0}$ in the same norm, since the difference of two
integrands is bounded by
$C e^{2\mu s}\|u-\widetilde u\|_\mu$.
Thus the map is a contraction for sufficiently negative $\tau_0$.
Differentiating its integral equation gives
$u,u',u''=O(e^{\mu\tau})$, and the differential equation gives
smoothness and the corresponding higher derivative estimates.

The constraint is preserved.  Indeed, the derivative of
\[
 \mathcal C=(S')^2-\sum_i(\alpha_i')^2-q^2
                  -\mathcal B_1-\mathcal B_2
\]
vanishes, because $\sum_i(v_r)_i=1$ and
$\mathcal B_r'=2(w_r\cdot\alpha')\mathcal B_r$.
Its limit as $\tau\to-\infty$ is zero by the Kasner relations, so
$\mathcal C=0$.  Hence the constructed functions solve all the
Einstein-scalar equations.  Since
$S=\tau+O(e^{\mu\tau})$, integration of the lapse from $-\infty$
gives $t=e^\tau(1+O(e^{\mu\tau}))$, and proves
\eqref{eq:einstein-optimal-separation-data} in proper time.

For the filling estimates, the central directions are $z_1,z_2$
dual to $\omega^1,\omega^4$.  On a sufficiently early tail the Schur
forms are uniformly comparable with
\begin{equation}
 H_t^0=\operatorname{diag}
    (t^{-1/10},t^{7/10},t^{7/10},t^{7/10}),
 \qquad
 A_t^0=\operatorname{diag}(t^{-4/5},t^{4/5}),
 \qquad L_t=0.
 \label{eq:einstein-optimal-schur-powers}
\end{equation}
The two nonzero brackets have exponents $3/10$ and $7/10$.
The central-component bound
\eqref{eq:anisotropic-bracket-matrix-bound} and its comparison
estimate therefore give
\[
 \kappa_{\rm opt}(T)\le\kappa_{\rm cen}(T)
       \le C\bigl(T^{3/5}+T^{7/5}\bigr)^{1/2}
       \le C' T^{3/10}.
\]
This order is attained.  Place a unit-area smooth ellipse in the
first horizontal pair on $[T/4,T/2]$, with coordinate amplitudes
$T^{(p_3-p_2)/2}$ and $T^{(p_2-p_3)/2}$.
For the central covector dual to $z_1$, its numerator in
\eqref{eq:kappa-opt} is one and its denominator is at most
\[
 C T^{p_2+p_3-p_1-1}=C T^{-3/10}.
\]
Here the dual central norm is $a_1^{-1}$, and the two amplitudes
balance the horizontal energy terms.  Taking this loop and covector
in the variational supremum proves the lower bound
$\kappa_{\rm opt}(T)\ge cT^{3/10}$.

To prove the second assertion in
\eqref{eq:einstein-optimal-separation}, fix $T>0$ and a compact
interval $J$ in a sufficiently early part of $(0,T)$.
Choose smooth planar based loops of signed area one.  For small
$s>0$ with $[s,2s]$ preceding $J$, place such a loop in the second
horizontal pair on $[s,2s]$, using the rescaled parameter $t/s$.
Place a loop of signed area $\varepsilon$ in the first pair on $J$
by multiplying a fixed loop by $\sqrt\varepsilon$.  Set both loops
equal to zero elsewhere.  Their direct sum is a based loop $\xi$
on a compact interval in $(0,T)$, with bracket area
\[
 Q_I(\xi)=\varepsilon z_1+z_2.
\]
On $J$, the central norm of this area has a positive lower bound
independent of $\varepsilon$, because its $z_2$ component equals one.
The contribution to $J_I(\xi)$ on $J$ is therefore at most
$C\varepsilon$.  On $[s,2s]$, the horizontal energy is at most
$C s^{-3/10}$, whereas
\[
 \sqrt{A_t(Q_I(\xi),Q_I(\xi))}
       \ge c\varepsilon s^{-2/5}.
\]
The early contribution is consequently at most
$C\varepsilon^{-1}s^{1/10}$.  Taking
$\varepsilon=s^{1/20}$ proves
\[
 J_I(\xi)\le C s^{1/20}\longrightarrow0.
\]
Thus $\kappa_{\rm an}(T)=\infty$ for every such $T$, and by
monotonicity for every $T$ in the proper-time domain.

All spatial Kasner exponents are less than one, so the cone integral
is finite.  The horizontal exponent $p_2=-1/20$ gives unbounded
diameter on every fixed open set of the cover.  Theorem
\ref{thm:two-step-obstruction} excludes every past continuous extension
on every discrete quotient.  The expansion is positive near the past
end.  The maximal homogeneous continuation has nonpositive spatial
scalar curvature and six spatial dimensions, so
Theorem~\ref{thm:homogeneous-future} gives future timelike and null
completeness.  This proves global $C^0$-inextendibility.

For the vacuum datum \eqref{eq:vacuum-optimal-separation-data},
\[
 \sum_i p_i^{\rm vac}=1,\qquad
 \sum_i(p_i^{\rm vac})^2
   =\frac{244+2\cdot78}{400}=1.
\]
The central exponents, the second horizontal pair and the sum of
the first horizontal pair are unchanged.  Thus both curvature
exponents remain $3/10$ and $7/10$, and the same Volterra construction
with constant scalar gives a vacuum solution with error $O(t^{3/5})$.
The central-component estimate still gives
$\kappa_{\rm opt}(T)=O(T^{3/10})$, and the balanced first-pair
ellipse gives the same lower bound because
$p_2^{\rm vac}+p_3^{\rm vac}=3/10$.
The early loop in the second pair
is unchanged, while the loop in the first pair remains on a fixed
compact time interval and still costs $O(\varepsilon)$.
Consequently the proof of $\kappa_{\rm an}=\infty$ is unchanged.
All vacuum exponents are less than one, and
$(3-\sqrt{78})/20<0$ is horizontal.  The same cone, past obstruction
and future completeness arguments apply.

Finally, all inequalities used above are strict.  In particular, if
$c_1,c_2$ denote the two central exponents and $r_1,r_2$ the second
pair's horizontal exponents, then
\[
 1+c_2-r_1-r_2>0,
 \qquad \delta:=r_1+r_2-1-c_1>0.
\]
For nearby unequal $r_1,r_2$, use on $[s,2s]$ a unit-area ellipse
whose two coordinate amplitudes are
$s^{(r_2-r_1)/2}$ and $s^{(r_1-r_2)/2}$.
Its energy is $O(s^{r_1+r_2-1})$, and the same argument yields
$J_I\le C(\varepsilon+\varepsilon^{-1}s^\delta)$.
The choice $\varepsilon=s^{\delta/2}$ again gives divergence of
$\kappa_{\rm an}$.  The Volterra construction, cone estimate and
negative horizontal exponent also persist in a sufficiently small
relative neighborhood of either displayed Kasner datum on the
appropriate vacuum or scalar Kasner constraint set.  This is a
statement about the diagonal
asymptotic family, with no assertion of genericity in the full space
of invariant or inhomogeneous initial data.
\end{proof}

\end{document}
